% Independent full-draft audit passed; see notes/de_A46_audit.md.
\subsection{A residue-prime normalizer criterion from zero p-rank}

\begin{statement}{A.46}{Proved}{Zero p-rank excludes noncuspidal pro-p normalizers}
Let \(Y\) be an affine hyperbolic curve over a number field \(L\),
let \(C\) be its smooth proper compactification, and put
\(J=\operatorname{Jac}(C)\). Suppose that at some place of \(L\)
above a rational prime \(p\), the abelian variety \(J\) has
potentially good reduction with geometric special \(p\)-rank zero.
In the geometrically pro-\(p\) arithmetic fundamental group of \(Y\),
every nontrivial closed procyclic subgroup of the geometric kernel
whose arithmetic normalizer has open image in \(G_L\) is contained
in a conjugate of a cusp inertia group.

The representation-theoretic assertion used here is stronger:
\(V_pJ=T_pJ\otimes_{\mathbf Z_p}\mathbf Q_p\) has no
one-dimensional \(\mathbf Q_p\)-subquotient after restriction to
\(G_{L'}\), for any finite extension \(L'/L\).
Both assertions include \(p=2\). No splitting of the arithmetic
normalizer is assumed.
\end{statement}

\begin{proof}
We first prove the representation assertion locally. Enlarge the
completion to a finite extension \(k/\mathbf Q_p\) over which \(J\)
has good reduction, and let \(\widetilde J\) be its special fibre.
Write \(k_0=W(\kappa)[1/p]\) for the maximal unramified subfield,
and \(f=[k_0:\mathbf Q_p]\). The covariant crystalline functor
on cohomology identifies, as Frobenius modules,
\[
 D_{\mathrm{cris},k}\bigl(H^1_{\acute{\mathrm e}\mathrm t}
              (J_{\bar k},\mathbf Q_p)\bigr)
       \simeq H^1_{\mathrm{cris}}(\widetilde J/W(\kappa))[1/p].
                                                        \tag{A.46.1}
\]
Here the good-reduction comparison is applied with constant
coefficients to the abelian scheme: in
\cite[Section 5, Corollary 4 and the subsequent Remarks 1--5,
printed pp.256--260]{RootFaltings2002}, invert the crystalline period
in the comparison and its inverse, and take Galois invariants.
There is no boundary divisor or logarithmic monodromy in this
smooth proper case. Equivalently this is the usual crystalline
comparison for the first cohomology of an abelian scheme.

The Newton slopes on the right of (A.46.1) lie in \([0,1]\).
The multiplicity of slope zero is the special \(p\)-rank, and
polarization pairs every slope \(a\) with \(1-a\).
Thus all slopes lie strictly between zero and one. These facts,
including invariance under finite residue extension, are recalled
in \cite[Sections 2.2.1 and 2.2.3]{RootAchterPries2013}.

Crystalline representations are closed under subquotients, and
\(D_{\mathrm{cris},k}\) is exact and preserves dimension;
see \cite[Theorem 5.2.1(1)--(3) and Proposition 9.1.6]{fghBrinonConrad}.
Suppose a one-dimensional subquotient existed. Duality converts
one of \(V_pJ\) into one of its dual
\(H^1_{\acute{\mathrm e}\mathrm t}(J_{\bar k},\mathbf Q_p)\).
Its crystalline module \(D\) is a rank-one \(k_0\)-isocrystal.
Choose a basis \(e\), and write
\[
       \varphi(e)=a e,\quad a\in k_0^\times,
       \qquad
       \varphi^f(e)=
         \left(\prod_{j=0}^{f-1}\sigma^j(a)\right)e .
                                                        \tag{A.46.2}
\]
For the normalization \(v_p(p)=1\), its unique Newton slope is
\[
 \frac1f v_p\left(\prod_{j=0}^{f-1}\sigma^j(a)\right)
                         =v_p(a)\in\mathbf Z,          \tag{A.46.3}
\]
because \(k_0/\mathbf Q_p\) is unramified. This remains true
however ramified \(k/\mathbf Q_p\) is. Exactness makes \(D\)
a subquotient of the isocrystal in (A.46.1). Its slope must
therefore be one of that isocrystal's slopes: after taking
\(\varphi^f\), characteristic polynomials multiply in an exact
sequence. This contradicts the strict interval \((0,1)\).
No extension of the coefficient field \(\mathbf Q_p\) has been
made; a line only after such an extension would not be excluded
by (A.46.3).

For a finite global extension \(L'/L\), choose a completion above
the indicated place and then a further finite local extension
giving good reduction. A \(\mathbf Q_p\)-line or subquotient for
\(G_{L'}\) remains one on this local group. Good reduction and
special \(p\)-rank zero persist under these extensions. The local
contradiction therefore proves the asserted global statement
for every \(L'\).

Finally let \(A\subset\pi_1(Y_{\bar L})^{(p)}\) have the
normalizer in the statement. If it were noncuspidal, A.11 would
give a nonzero image of \(A\) in \(V_pJ\), whose rational span
is a one-dimensional \(\mathbf Q_p\)-subrepresentation of an open
arithmetic subgroup. More explicitly, geometric inner actions
are trivial on proper homology, so the normalizer's action on
this line factors through its open image in \(G_L\).
This is precisely a line over a finite extension excluded above.
The argument uses the actual proper-homology map and requires
neither a section of the normalizer nor any lifting of fixed
vectors from a rational completion.
\end{proof}
